\chapter{绪论}

\section{研究背景与意义}
\label{sec:chap1-background}

\subsection{多模态大语言模型的训练}

深度神经网络（Deep Neural Network, DNN）在计算机视觉（Computer Vision, CV）、自然语言处理（Natural Language Processing, NLP）、语音识别（Speech Recognition, SR）等众多前沿领域取得了显著成效，近年来逐步成为众多关键前沿技术的核心基础\cite{galvatron2022}。
2017 年，Vaswani 等人提出基于自注意力机制（Self-Attention）的 Transformer 架构%
\cite{vaswani2017}，
摒弃了传统的循环神经网络（Recurrent Neural Network, RNN）与卷积神经网络（Convolutional Neural Network, CNN）结构，由于其适合并行计算且擅长捕捉长上下文信息，Transformer 架构迅速成为自然语言处理与跨模态建模等领域的主流模型结构。

在自然语言处理领域，基于 Transformer 的预训练大语言模型（Large Language Models, LLMs）持续刷新多项基准测试的成绩，并推动模型能力由特定任务优化向通用知识学习演进。
以 BERT\cite{bert} 为代表的带掩码的双向编码预训练范式在语言理解任务中取得了突出效果，与此同时，GPT-2\cite{gpt2}、T5\cite{t5} 所采用的自回归（Autoregressive）与编码器-解码器（Encoder-Decoder）机制在文本生成场景中展现出更强的能力。
随着模型规模进一步扩展，GPT-3\cite{brown2020gpt3} 及后续 GPT-4\cite{openai2023} 在少样本（Few-Shot）与零样本（Zero-Shot）条件下呈现出极强的语言表达与任务迁移能力。

在计算机视觉领域，Transformer 架构也取代了传统的卷积神经网络，成为视觉模型的主流选择。Vision Transformer（ViT）\cite{vit2020} 及其规模化扩展工作\cite{scaling_vit} 表明，图像经块划分后输入 Transformer 层进行特征提取，可在 ImageNet 等基准数据集上达到并在部分配置下超越传统卷积神经网络的性能。
由此可见，Transformer 架构具备跨模态扩展的能力，其适用范围并不局限于文本建模。

随着单模态模型的能力已经接近甚至超越人类水平，将语言、视觉、音频、表格等多模态信息纳入统一的模型框架进行联合表示与计算，成为当前学术界和工业界研究的热点，这种可以同时处理两种或两种以上模态信息的能力被称为多模态（Multimodality）。
在机器学习与人工智能中，常见模态包括文本、图像、视频、音频、表格以及各类传感器数据，这种多种模态的信息更接近人类接触到的复杂信息。
多模态大语言模型（Multimodal Large Language Models, MLLMs）是指在原有大语言模型的基础上，增加对图像、声音等其它模态的感知与生成能力的一类模型%
\cite{blip2,llava2023,flamingo2022}。
MLLMs 的先驱工作 CLIP\cite{clip2021}（Contrastive Language-Image Pre-Training）通过使用图像文本对进行训练，使模型拥有理解并描述图片的能力，DALL-E\cite{dalle2021} 则与 CLIP 相反，展示了从文本生成图像的可行性。

近年来，MLLMs 在架构与训练范式上发展迅速。
BLIP-2\cite{blip2} 提出 Q-Former 连接冻结的视觉编码器模型与 LLM 基座模型进行训练；
InstructBLIP\cite{instructblip2023} 在此基础上加强了视觉信息的提取能力，在多项零样本任务上取得 SOTA 结果；
LLaVA\cite{llava2023} 与 LLaVA-1.5\cite{llava15_2023} 采用简单线性投影层实现视觉与语言模态对齐，验证了轻量的连接器与大语言模型组合方案的有效性；
Qwen-VL\cite{qwenvl2023}、InternVL\cite{internvl2024} 等系列多模态大语言模型则进一步在长上下文、高分辨率与多图理解上扩展了模型能力。
这类模型既保留传统 LLMs 在语言理解、推理与生成上的强大能力，又能够直接处理图像、音频等多模态输入，实现跨模态交互与推理任务，因而被广泛认为是通往通用人工智能（Artificial General Intelligence, AGI）的最有前景的技术手段\cite{minigpt4,mllm_survey_acl2024}。

从功能表现来看，MLLMs 可完成图像描述、视觉问答、图文对话、文档理解，以及基于文本的图像或视频生成等复杂任务，在自动驾驶、智能助手等应用场景中具有极大应用价值。
在垂直领域，MLLMs 已延伸至医疗影像理解\cite{medical_vlm2024}等。
多模态大语言模型的相关评估常使用 VQAv2\cite{vqav2}、GQA\cite{gqa2019}、MMMU\cite{mmmu_benchmark2024} 等多模态基准，不断更新的基准测试也对模型规模与训练数据提出了更高要求。
从结构上看，MLLMs 通常包含以下核心组件，如图~\ref{fig:mllm_components} 所示：
（1）Modality Encoder，负责将图像、视频等非文本输入编码为特征表示，常采用预训练的视觉模型（如 CLIP、ViT）；
（2）Input Projector，负责将编码后的多模态特征映射到与 LLM 词嵌入对齐的语义空间，常见形式包括线性层、Q-Former\cite{blip2} 等轻量级模块；
（3）LLM Backbone，即大语言模型主体，承担主要的内容理解与生成任务，处于 MLLMs 架构的核心地位，多模态能力通常是在其基础上的扩展与增强；
（4）Output Projector（在需要多模态生成时），将 LLM 输出映射到与目标模态生成器对齐的表示空间；
（5）Modality Generator（在需要多模态生成时），负责生成图像、音频等目标模态，常采用扩散模型（如 Stable Diffusion）等。

\begin{figure}[htbp]
  \centering
  \includegraphics[width=\linewidth]{mllm_components.pdf}
  \caption{多模态大语言模型的核心组成部分}
  \label{fig:mllm_components}
\end{figure}

多模态大语言模型能力的拓展既依赖上述编码器、投影层、基座模型与可选生成器等模块的协同设计，也依赖参数与数据规模的提升。
参数和数据规模决定了模型可达的效果上限，而多模块的协同设计则影响模型的训练过程。
在大语言模型占据主体参数量的架构下，针对 LLMs 总结出的扩展定律（Scaling Laws）与数据规模规律，同样深刻影响着 MLLMs 训练的资源需求。
已有实证研究表明，模型性能通常会随参数规模与训练数据规模的增加而稳定提升，这一现象被总结为“扩展定律”（Scaling Laws）\cite{kaplan2020scaling}。
为追求更强性能，工业界与学术界不断推出参数量更大、数据规模更大的模型，
从数亿参数（如 BERT）到千亿、万亿参数（如 GPT-3、Gopher\cite{gopher2022}、DeepSeek-V3\cite{deepseek2024}、Qwen2.5\cite{qwen2024}、
LLaMA 2/3\cite{llama2_2023,llama3}、Gemma\cite{gemma2024}、Mistral\cite{mistral2023} 等开源与闭源 LLMs），
单卡或单机 GPU 的显存与算力已无法容纳模型的完整训练流程%
\cite{megatron2021,megatron2020}。
相应地，多模态大语言模型的预训练与指令微调所需的数据量也急剧增长，从数百万图文对到数十亿级样本，单机的运算速度已经远远无法满足企业对模型迭代速度的需求。
不仅如此，长上下文与多模态联合基准\cite{long_context_benchmark2024} 的普及进一步推高了对训练规模与效率的需求。
在上述参数规模与数据规模的限制下，单个计算设备的显存与算力均难以独立支撑完整的训练流程，因此必须借助分布式并行训练技术，将模型参数、优化器状态、激活值与数据分散至多台设备，通过节点间的通信协同完成多模态大语言模型的训练%
\cite{zero2020,zero2021}。

目前，支撑大规模 MLLMs 训练的主流手段包括数据并行、张量并行、流水线并行、序列与专家并行等多维混合并行技术\cite{pytorch_ddp,megatron2020,megatron2021,gpipe,pipedream2019,gshard}，以及面向参数与优化器状态的 ZeRO 系列显存分片技术\cite{zero2020,zero2021}；上述策略常与混合精度训练\cite{mixed_precision2018}、Flash Attention\cite{flash_attention_2022,flashattention3_2024} 等相结合，共同支撑当前大规模模型训练任务。
在采用上述并行策略的同时，训练过程仍面临较大的显存压力。
由于前向传播中需要保存中间激活以供反向传播时计算梯度使用，中间激活的显存占用随序列长度近似指数增长，在单设备上常构成主要显存压力，从而限制模型训练的上下文长度，进而限制模型能力上限。
ZeRO 等分片技术主要缓解参数与优化器状态的显存占用，对单设备上激活占用的降低有限。
为此，重计算（亦称梯度检查点）通过在前向传播中丢弃中间激活、在反向传播时重算，以额外前向计算换取显存空间，使模型在相同硬件条件下能够采用更大批次或更长序列进行训练\cite{chen2016checkpoint,megatron2023}。
由于如今的模型训练上下文迅速增长，重计算成为当前大规模训练的显存优化关键技术，因此 Megatron-LM、DeepSpeed 等主流并行训练框架均支持重计算技术。

然而，上述策略在设计时多假设集群同构且模型是单一 Transformer 层堆叠的语言模型，其划分粒度、负载均衡与重计算策略均围绕“各模型层计算与显存相对均匀”这一前提展开。
对于 MLLMs 而言，会引入两个额外问题：其一是模型结构异构，MLLMs 各子模块在算子类型、参数量与激活体积上差异显著。
编码器以卷积或 ViT 块为主，投影层参数量较小，而 LLM 基座由大量同构 Transformer 层堆叠且占据绝大部分参数。这种异构性使得传统基于“按层均分”假设的模型划分方案难以适用。
其二是分阶段冻结导致的动态负载变化。
当某一模块被冻结时，其反向传播计算开销与优化器状态的显存开销显著下降，各模块的计算与显存占比随阶段切换而动态变化。
例如，在冻结编码器阶段，编码器侧的显存与计算负载骤减，而 LLM 基座承担大部分梯度计算，切换到冻结 LLM 阶段后，负载分布又会发生变化。
这种动态变化特性使得固定的模型划分配置难以在全训练流程中保持高效。

\subsection{异构 GPU 集群下的并行训练}

多模态大语言模型的训练效率与成本高度依赖底层硬件平台。
目前，工业界与学术界普遍采用 NVIDIA GPU 作为 MLLMs 训练的主要加速器，并通过多卡、多机集群进行分布式训练。
人工智能的快速发展也推动了硬件迭代的速度，大约每 1--2 年，NVIDIA 就会推出使用新一代计算架构的多款不同定位的 GPU，这给 MLLMs 的训练带来了新的挑战。

图~\ref{fig:nvidia_gpu_evolution} 描述了近年来 NVIDIA GPU 的架构演进。
2016--2017 年的 Pascal 计算架构（如 NVIDIA P100）与 Volta 计算架构（NVIDIA V100）奠定了数据中心 GPU 与 Tensor Core 的基础。
2018 年 Turing 架构（RTX 20 系）引入消费级光追与 INT8 推理。
2020 年 Ampere 架构推出 A100（40GB/80GB HBM2e）与消费级 RTX 30 系列（如 RTX 3090 24GB），支持 PCIe 4.0 与 NVLink 3.0，成为当前许多实验室与云平台的主力\cite{nvidia_a100}。
2022 年搭载 Hopper 架构的 H100（80GB/96GB HBM3）针对大语言模型训练做了专门优化\cite{nvidia_h100}，并广泛配备 NVLink 4.0 与 NVSwitch，单机多卡带宽可达数百 GB/s，
同年 Ada Lovelace 架构的 RTX 40 系列（如 RTX 4090 24GB）在消费级市场提供高性价比算力。
2024 年起，Blackwell 架构的 B100/B200 等数据中心 GPU 逐步量产，采用 HBM3e、更高带宽第五代 NVLink 与 PCIe 6.0，进一步拉大与旧代设备的性能与互联差距。
据 NVIDIA 在 GTC 等场合公开发布的产品路线图\cite{nvidia_blackwell}，后续还将推出 Rubin 架构（约 2026 年）、Rubin Ultra（约 2027 年），以及面向推理与能效深度优化的 Feynman 架构（约 2028 年）。
与此同时，不同型号的 GPU 在 PCIe 版本（3.0/4.0/5.0/6.0）、是否配备 NVLink/NVSwitch、以及 InfiniBand 与 RoCE 等节点间网络上的不同，也会带来通信带宽的差异。

\begin{figure}[htbp]
  \centering
  \includegraphics[width=\linewidth]{nvidia_gpu_evolution.pdf}
  \caption{NVIDIA GPU 架构与代表性产品演进}
  \label{fig:nvidia_gpu_evolution}
\end{figure}

在大多数高校、研究所与中小企业的实际环境中，受预算与采购周期限制，GPU 集群往往经历多次分批采购，而旧设备报废周期较长，集群中常同时包含多代 GPU（例如部分 A100 或 H100 与大量 RTX 3090/4090、甚至更早的 V100/RTX 2080），单机内也可能混用不同显存与互联配置，难以做到“全集群同构”，由此形成\emph{显存异构}与\emph{算力异构}并存的局面\cite{whale,hetpipe2020}。
然而，主流并行训练框架（如 Megatron-LM、DeepSpeed）在设计时多假设同构设备与统一拓扑\cite{megatron2021,zero2020}，在异构集群上直接使用容易出现以下问题：负载无法按算力与显存合理分配，低显存设备成为瓶颈而高端卡显存资源利用率低，以及计算最慢的节点拖慢整体迭代速度\cite{whale,hetpipe2020}。
对于单层计算与显存相对均匀的传统大语言模型，尚可通过“按层数或算力比例”划分在一定程度上缓解异构带来的影响。
而对于 MLLMs，其模块间在层类型、参数量与激活体积上的固有差异与设备能力差异相互叠加，使得简单的模型划分难以实现负载均衡。
因此，在显存与算力异构的 GPU 集群上高效训练 MLLMs，是学术界和工业界共同面临且亟待解决的问题。

从系统层面看，异构带来的问题不仅在于单卡能力差异，还在于拓扑与调度。
节点内 NVLink 与节点间 InfiniBand 的带宽可相差一个数量级以上，若流水线并行或张量并行的划分未考虑链路差异，跨节点通信极易成为瓶颈。
同时，MLLMs 各训练阶段对编码器、投影层与 LLM 基座的冻结策略不同，导致同一模型在不同训练阶段下各模块的计算与显存开销发生明显变化。
若沿用固定的模型划分与调度，难以在整轮训练中保持负载均衡。
即便划分与调度已得到优化，在异构流水级上采用全局统一的重计算配置仍难以同时满足各设备的显存约束。
因此，面向异构集群的 MLLMs 训练需要同时考虑硬件拓扑、模型结构、多训练阶段以及流水级差异化的显存—计算权衡，通过自适应的划分、高效的调度与按流水级配置的重计算策略协同提升训练效率。

综上所述，硬件层面的设备异构与模型层面的模块异构、分阶段冻结相互交织，使得异构集群上的 MLLMs 训练同时面临负载分配、算存资源利用、显存—计算权衡等多维约束，构成了当前有限资源下高效训练的核心难点。
随着云上与混合部署中弹性训练与按需资源供给的普及，集群内设备型号与拓扑将更加多样，这对异构感知的划分、调度与重计算配置提出了更高要求\cite{frenzy_memory2024,nvidia_training_guide,arnold2025}。
为此，本文基于 Megatron-LM 主流并行训练框架，针对 MLLMs 的模块结构异构与分阶段冻结特性，结合异构 GPU 集群中显存与算力的多维差异，研究细粒度的混合并行划分与重计算优化技术，旨在提出一套面向异构 GPU 环境的自适应并行训练方案，以提升异构 GPU 集群下 MLLMs 的训练效率。

\section{国内外研究现状与挑战}

本节从文献角度归纳与本文密切相关的研究脉络，先概述多模态大语言模型并行训练策略研究现状及其同构假设的局限，再综述面向异构 GPU 集群的并行训练工作，最后归纳重计算相关研究及其在异构 GPU 场景下的不足。

\subsection{多模态大语言模型并行训练策略研究现状}

Kaplan 等\cite{kaplan2020scaling} 通过系统实验提出扩展定律（Scaling Law），指出模型性能会随参数量与训练数据量的增加而提升；Chinchilla\cite{chinchilla2022} 进一步从计算最优角度表明参数规模与训练词元（token）数应近似同比例扩展，后续研究\cite{beyond_chinchilla2024} 又将该缩放规律拓展到推理成本约束下的最优配置。
基于上述研究，多模态大语言模型普遍沿着“更大参数规模+更大数据规模”的路径持续演进。

自 Transformer\cite{vaswani2017} 提出以来，语言模型规模呈持续扩张趋势，参数量由 BERT\cite{bert} 的亿级逐步发展到 GPT-2\cite{gpt2}、T5\cite{t5} 的十亿级，并进一步扩展至 GPT-3\cite{brown2020gpt3}、GPT-4\cite{openai2023} 以及 DeepSeek-V3\cite{deepseek2024}、Qwen2.5\cite{qwen2024} 等模型所代表的千亿级规模。
在此基础上，MLLMs 通过引入视觉编码器并与语言模型进行跨模态对齐，形成了以 BLIP-2\cite{blip2}、InstructBLIP\cite{instructblip2023}、LLaVA/LLaVA-1.5\cite{llava2023,llava15_2023}、Flamingo\cite{flamingo2022}、Qwen-VL\cite{qwenvl2023}、InternVL\cite{internvl2024} 和 MiniGPT-4\cite{minigpt4} 等为代表的技术谱系，其模型规模已覆盖数十亿至数千亿参数，训练数据规模也由百万级图文对扩展至十亿级样本（如 LAION-5B\cite{laion5b2022}）。
相关工作已从模型架构、跨模态对齐策略、数据构建与系统效率等维度对该领域进行了系统总结\cite{mllm_survey_acl2024,comprehensive_mllm_survey2024,mllm_datacentric2024}。
单卡或单机在显存容量和计算吞吐上已难以支撑完整训练流程\cite{megatron2021,megatron2020}，并行与分布式训练因此成为大模型与 MLLMs 训练的基础技术路径\cite{zero2020,zero2021}。

除了模型的参数规模，训练数据的规模同样持续扩大。图像方面，Open Images\cite{open_images}、COCO\cite{coco2014} 等海量图片数据集需 TB 级存储，LAION-5B 等开放图文对达数十亿样本。
文本方面，Common Crawl、The Pile\cite{pile2020} 等为 LLM 预训练提供数百 GB 至数万亿词元量级的语料。
视频方面，YouTube-8M\cite{youtube8m}、WebVid\cite{webvid2021} 等基准需 PB 级存储与高吞吐数据管线，多模态理解与推理则依赖 VQAv2\cite{vqav2}、TextVQA\cite{textvqa2019} 等数据集与长上下文基准\cite{long_context_benchmark2024}。
在单机无法存储完整数据集与模型的前提下，并行训练成为必然选择。

在并行策略方面，工业界与学术界在同构 GPU 集群上已形成以模型与数据切分为核心的多维混合并行，以及以状态分片为核心的 ZeRO 系列两大技术路线，并被 Megatron-LM、DeepSpeed 等并行训练框架广泛应用于千亿级参数 MLLMs 训练\cite{megatron2021,zero2020,zero2021}。
多维混合并行通过在不同维度上划分模型或数据以降低单卡负载：\emph{数据并行}\cite{pytorch_ddp,sergeev2018horovod} 在各设备组上复制完整模型、按批次划分样本，通信集中在梯度 AllReduce，实现简单且扩展性好；
\emph{张量并行}\cite{megatron2020,megatron2021} 将单层内矩阵按行或列切分到多卡，降低单卡显存，是千亿级稠密模型的主流选择；
\emph{流水线并行}将模型按层划分为若干流水级并分配到不同设备，有效支持超大规模模型的并行训练，但需借助更精细的调度以缓解流水线气泡、提高 GPU 利用率。
此外，\emph{序列并行}\cite{megatron2021}、\emph{专家并行}\cite{gshard} 等可与张量、流水线并行组合，以进一步扩展训练规模。
NVIDIA 的 Megatron-LM\cite{megatron2021,megatron2020} 联合使用张量、流水线与数据并行，已支撑在数千张同构 GPU 上训练千亿至万亿级参数的大语言模型。
零冗余优化器 ZeRO\cite{zero2020,zero2021} 在数据并行的基础上将优化器状态、梯度与参数分片存于各卡、按需聚合，显著降低单卡显存，但引入较多卡间通信；ZeRO-Infinity\cite{zero2021} 进一步将部分状态卸载至 CPU 或 NVMe；ZeRO++\cite{zeroplusplus} 则通过量化与通信优化削减跨节点开销。

% 在同构集群的自动并行方向，Galvatron\cite{galvatron2022} 针对 Transformer 将策略搜索与编译优化结合；Aceso\cite{aceso2024} 通过迭代瓶颈缓解联合优化通信与显存；UniAP\cite{uniap2025} 将层间与层内并行统一为混合整数二次规划以搜索全局配置；Colossal-Auto\cite{colossal_auto} 则将自动并行与重计算联合自动化。

% 上述方法多在\emph{同构}集群上为全部设备配置统一的重计算粒度与子模块集合。

然而，主流训练框架在设计时多假设\emph{同构} GPU 集群与相同 Transformer 层堆叠的语言模型，对 MLLMs 的\emph{模块异构}（编码器、投影层、LLM、生成器结构各异）与\emph{分阶段冻结训练}考虑不足。
以 Megatron-LM 为例，其在同构集群上训练 MLLMs 时，编码器与生成器的并行方式通常被动跟随 LLM 基座，缺乏面向各子模块算力与显存需求的细粒度映射，典型策略如图~\ref{fig:megatron_mllm} 所示\cite{disttrain2025}。
面向 MLLMs 的模块结构，DistMM\cite{distmm2024} 针对子模块差异组合张量并行与数据并行，DistTrain\cite{disttrain2025} 通过解耦模态编码器、LLM 基座与模态生成器的训练流程分别分配资源且配置并行策略，在同构集群上提高了训练 MLLMs 时集群的利用率。
在显存优化方面，Korthikanti 等\cite{megatron2023} 在 Megatron-LM 中提出序列并行与选择重计算，与张量并行协同以降低激活显存；Colossal-Auto\cite{colossal_auto}、Mist\cite{mist2025}、AdaPipe\cite{adapipe2024} 等工作进一步将重计算与自动并行或流水线划分联合搜索。

总体而言，上述面向 MLLMs 的并行训练方法多建立在同构 GPU 集群假设之上，忽视了异构设备间显存容量与计算吞吐的差异，对\emph{分阶段冻结训练}所引起的计算与显存负载的动态变化也缺乏针对性适配。
即便在同构集群中沿用上述策略，仍易诱发流水级负载失衡与局部显存瓶颈；在设备能力分布相对均衡而子模块算存需求差异显著的条件下，上述问题将进一步加剧。

\begin{figure}[htbp]
  \centering
  \includegraphics[width=\linewidth]{fig02_megatron_mllm.jpg}
  \caption{使用 Megatron-LM 框架在同构集群上训练多模态大语言模型的典型并行配置}
  \label{fig:megatron_mllm}
\end{figure}

\subsection{异构 GPU 集群并行训练策略研究现状}

受预算与采购周期限制，实际训练集群往往由多代、多型号 GPU 混合组成，呈现显著的显存与算力异构\cite{whale,hetpipe2020}。
将面向同构集群的并行技术与重计算策略直接用于异构 GPU 集群时，不同显存容量、算力与互联带宽的 GPU 设备会被视为同构的设备，低显存或低算力节点易成为整个系统的瓶颈，高端设备资源利用率不足，整体训练效率下降\cite{whale,hetpipe2020}。
针对上述问题，学术界已开展异构 GPU 感知的并行规划与资源调度研究。

早期工作如 HetPipe\cite{hetpipe2020} 将流水线并行与数据并行结合，借助虚拟工作节点（VW）在异构设备间分配负载，但未显式建模设备排列与通信拓扑，也未针对 MLLMs 的模块异构与分阶段冻结开展联合设计。
Whale\cite{whale} 从显存与算力差异出发进行硬件感知负载均衡，按异构 GPU 算力比例划分模型与数据批次，但其划分主要围绕设备算力展开，对模型模块异构带来的计算与显存差异覆盖仍不充分。
AMP\cite{amp2022} 在异构环境下自动搜索并行策略，综合考虑模型异构与设备异构，通过开销模型与动态规划得到近似最优配置；然而其划分粒度偏粗，通信与显存估计精度不足，难以支撑更细粒度的高效映射。

近年来，面向异构 GPU 集群的分布式训练系统在“异构感知并行规划—资源编排—代价建模”等方向上持续演进：HetHub\cite{hethub2024} 统一规划通信器、性能预测器与自动并行以适配多厂商 GPU；HexiScale\cite{hexiscale2024} 将多种并行划分形式化，并通过层次图划分在异构环境中分配计算；Espresso\cite{espresso2025} 联合优化异构资源供给与流水级放置；Metis\cite{metis2024} 以异构感知搜索实现快速自动并行规划；HAP\cite{hap2024} 将 SPMD 并行优化建模为程序合成问题，联合搜索张量分片与通信方式。
总体而言，上述异构感知工作多数仍偏向通用语言模型，且仅对\emph{设备异构}做单一维度优化，对于 MLLMs 特有的\emph{模块异构}、\emph{分阶段冻结}等特性则缺乏针对性调优。

在异构 GPU 集群上，面向同构集群的全局统一重计算配置适用性也进一步下降：显存不足的流水级若配置偏保守，仍可能因中间激活过大而 OOM，导致训练过程无法进行；显存充裕的流水级则可能被强制对齐同等强度的重计算，引入冗余前向计算并延长迭代时间。
在 MLLMs 模块异构与分阶段冻结并存的场景下，如何在与划分、调度协同的前提下按流水级差异化配置重计算，仍缺乏系统化研究。

除多维混合并行外，以 ZeRO 为代表的状态分片路线在异构 GPU 集群上训练 MLLMs 时亦面临局限：ZeRO\cite{zero2020,zero2021} 引入大量 All-Gather/All-Reduce 通信，在多节点集群中易成为瓶颈；MiCS\cite{mics2022} 难以结合设备算存差异与冻结阶段切换实现细粒度、动态的负载自适应；ZeRO++\cite{zeroplusplus} 的量化与通信优化可能带来精度与通用性代价；AMSP\cite{amsp} 等侧重跨节点通信削减，对集群内存储、算力与带宽的细粒度异构利用仍有限。
因此，在异构 GPU 集群上高效训练 MLLMs，仍需在 Megatron-LM、DeepSpeed 等现有框架基础上，引入对设备能力、模型模块结构与训练阶段的联合感知，实现划分、调度与重计算的协同自适应优化。

\subsection{研究问题与挑战}

围绕异构 GPU 集群上的多模态大语言模型训练，本文聚焦的核心研究问题是：在设备显存、算力与带宽不一致的真实异构 GPU 集群环境下，如何针对 MLLMs 的模块异构结构与分阶段冻结训练的特性，联合优化模型划分、流水线调度与重计算配置，在满足显存约束与收敛正确性的前提下最小化迭代时间，提升端到端训练速度。

为回答这一问题，本文面临以下三项关键挑战。
第一，模块异构与设备异构耦合导致负载失衡。MLLMs 由编码器、投影层与 LLM Backbone 等多个模块构成，不同模块在计算量与显存占用上差异显著。
当其与异构 GPU 的算存能力差异叠加时，传统按层均分或按单一硬件指标划分模型的方法容易形成瓶颈流水级，导致整体吞吐受限。
第二，冻结阶段切换使最优并行策略动态变化。
在不同冻结策略中，可训练参数集合动态变化，反向计算与梯度同步开销随之动态变化，静态划分与固定调度难以在全流程保持高效。
第三，统一的重计算策略难以兼顾显存安全与训练效率。在使用更大批次与更长序列训练时，部分流水级存在显存瓶颈并可能发生 OOM。

基于上述挑战，本文的研究目标是构建一套面向异构 GPU 集群上多模态大语言模型训练的联合优化框架：在模型划分层面考虑异构 GPU 设备与冻结导致的模型算存需求的变化，在调度层面利用冻结信息剪枝冗余反向与通信，在重计算策略层面按流水级配置差异化重计算策略，从而形成“划分—调度—重计算”多级优化框架。
具体而言，挑战一主要对应框架的划分层，侧重模块异构与设备异构耦合下的负载均衡；挑战二则要求划分与调度随冻结状态协同调整，并消除冻结场景下的冗余反向与通信，对应框架的划分层与调度层，由第三章 MMoH 的并行策略规划器与提前终止调度器共同解决；挑战三主要对应框架的重计算层，由第四章 HR-MMoH 在 MMoH 基础之上按流水级进行差异化重计算配置，以支撑更大批次与更长序列训练并进一步压缩迭代时间。

\section{研究内容与贡献}

围绕异构 GPU 集群上多模态大语言模型的高效训练这一总体目标，本文将研究工作分解为两个层次递进、相互衔接的子问题。
第一，划分与调度子问题：在设备显存与算力异构、且训练过程伴随分阶段冻结切换的条件下，如何自动生成面向多模态大语言模型的细粒度模型划分、设备映射与流水线调度方案，使各流水级计算负载与异构设备能力相匹配；如何随各训练阶段的冻结状态生成相应的划分与映射配置，并在流水线执行中剪枝连续冻结流水级的冗余反向计算与级间通信，以提升 GPU 利用率与训练吞吐。
第二，重计算子问题：在第一子问题所确定的并行执行骨架（设备拓扑、模型划分与冻结感知调度）之上，当训练规模向更大批次与更长序列扩展时，如何按流水级配置差异化的重计算策略，在各设备显存约束下尽可能降低重计算引入的额外计算开销，进一步缩短迭代时间。

针对上述两个子问题，本文依次提出 MMoH 与 HR-MMoH 两种方法，分别针对划分—调度与重计算两个层次的优化任务，研究内容与贡献如下。

\subsection{MMoH：多模态大语言模型在异构 GPU 集群上的混合并行方法}

针对异构 GPU 环境中显存与算力差异导致的设备利用率低、模块结构差异引发的模型划分困难，以及分阶段冻结训练造成的计算与显存需求动态变化等问题，本文提出一种面向异构 GPU 集群的多模态大语言模型混合并行训练方法。该方法通过为不同计算设备分配不同规模的模型分片，实现对计算资源与显存空间的协同利用，从而提升训练吞吐并缩短迭代时间。进一步地，为适应多阶段冻结训练过程，该方法能够感知各阶段冻结状态并与并行划分方案协同调整，同时在连续冻结的流水级上省略不必要的反向计算与级间通信，在保证模型收敛性的前提下进一步提升异构设备的有效利用率与整体训练效率。
MMoH 承担本文两级方法中的划分与调度层，其输出包括设备拓扑、模型映射及冻结感知的流水线调度策略，构成后续重计算优化的优化对象与约束条件。

\subsection{HR-MMoH：多模态大语言模型在异构 GPU 集群上的重计算方法}

针对现有重计算策略多采用全局统一配置、未考虑异构 GPU 集群中设备显存与算力差异的问题，本文提出一种面向异构设备与多模态大语言模型的重计算优化方法。
HR-MMoH 以 MMoH 的划分与调度结果为输入，将重计算决策粒度由全局统一下沉至流水级差异化：以各设备显存上限为约束，对显存受限流水级采用较强的重计算以降低激活驻留，对显存充裕流水级则减少冗余重算，从而在满足显存约束的前提下最小化重计算引入的额外计算开销。
HR-MMoH 不改变 MMoH 的划分与调度策略，而是在其基础之上与之共同构成“划分—调度—重计算”多级优化链路，进一步提升异构 GPU 集群上 MLLMs 的端到端训练效率。


\section{论文组织结构}

本文共分为以下五个章节，具体结构如下：

第一章为绪论。首先从多模态大语言模型的训练需求与异构 GPU 集群的现状两个方面阐述研究背景与意义，随后分别从多模态大语言模型并行训练和异构 GPU 集群并行训练两个维度梳理国内外研究现状，最后介绍本文的研究内容与主要贡献。

第二章为相关工作。首先介绍多模态大语言模型的模型架构、训练过程以及数据并行、模型并行与 ZeRO 等主流并行训练技术，其次综述异构 GPU 集群上的并行训练策略与工业实践，最后阐述重计算技术的基本原理、分类及主流实现的局限性。

第三章介绍多模态大语言模型在异构 GPU 集群上的高效混合并行技术 MMoH，重点阐述需求驱动的设备拓扑识别、复合粒度模型抽象、冻结感知划分及提前终止流水线调度方法，并给出实验评估。

第四章介绍多模态大语言模型在异构 GPU 集群上的重计算技术 HR-MMoH，详细说明如何在异构 GPU 集群与不同模型模块间实现重计算策略的细粒度自适应配置并进行实验验证。

第五章为总结与展望。概括本文主要创新点，归纳本文的研究局限，并对后续研究问题与思路作简要展望。
